| Topic               | SotA (theory, options, industry standards)                                                                                                           | Methods (what we actually did)                                                    |
| ------------------- | ---------------------------------------------------------------------------------------------------------------------------------------------------- | --------------------------------------------------------------------------------- |
| ML for sequences    | Brief overview of approaches (HMM, TF-IDF, classical, mention of deep/BERT as existing direction)                                                    | HMM + TF-IDF specifics, Python libs, Jupyter, our evaluation                      |
| Web architectures   | Monolith vs microservices, SPA vs MPA, client-server, polyglot persistence as pattern                                                                | Our containerized SPA + FastAPI + Celery stack                                    |
| Frontend principles | SPA, PWA, SSR vs CSR, component-based approach, Virtual DOM + JSX as concepts, state management paradigms, browser caching/storage                   | React+TS specifics, our hooks, our caching choice, our heatmap                    |
| Backend / API       | REST vs SOAP vs GraphQL, HTTP details (headers, cookies, CORS), middleware concept, ORM concept, auth standards (OAuth, JWT, sessions)               | FastAPI specifics, SQLModel, our CORS+middleware, our JWT+hashing impl            |
| Data storage        | SQL vs NoSQL types, in-memory DBs, file storage, object storage (S3), ML model storage                                                               | Our MySQL + MongoDB + Elasticsearch + Redis choice, how we store files and models |
| Deployment          | Dedicated vs cloud, SaaS, containers vs VMs, Docker vs K8s, CI/CD (GH Actions, GitLab CI, Jenkins), IaC (Terraform), envs (dev/uat/prod), serverless | Our Docker Compose + Nginx reverse proxy + optional HTTPS setup                   |
| UX / design         | UCD principles, chunking, personas, HTA (general)                                                                                                    | Figma process, our personas, our wireframes, testing with 4 users                 |
| AI in dev           | —                                                                                                                                                    | Our Codex/Claude Code/Mistral experiments                                         |

\subsection{Machine learning for sequential medical codes}
\label{subsec:ml_methods}

The main analytical task of this work is to process longitudinal electronic health records. Each patient history is represented as a sequence of five-digit medical procedure codes. These sequences vary a lot in length and the time intervals between visits are irregular. So traditional tabular machine learning methods cannot be applied directly. For this reason we explored two different approaches. The first one is probabilistic sequence modeling with Hidden Markov Models (HMM). The second one is vector space modeling with Term Frequency-Inverse Document Frequency (TF-IDF) combined with classical classifiers.

\subsubsection{Probabilistic sequence modeling (Hidden Markov Models)}
HMM is a probabilistic model for sequence data. It assumes that the observed sequence is generated by an underlying process which moves through a finite set of hidden states. At each step the process is in one of the hidden states and emits one observable output. The probability of moving from one hidden state to another is given by a transition matrix. The probability of emitting a specific output from a given state is given by an emission distribution~\cite{hmm_medical_time_series_2022}. In the healthcare setting, the hidden states represent the underlying stage of the patient's disease, for example healthy, early-stage, or advanced. These stages are not written in the records directly. The observable outputs are the discrete medical procedure codes recorded during hospital visits.

HMMs were chosen for initial experimentation for several reasons. They handle variable-length sequences naturally, so there is no need to pad or align the data. Another reason is that they can model temporal dependencies between visits without assuming regular time steps~\cite{hmm_bioinformatics_2025}. Also, HMMs are interpretable. The learned transition probabilities and emission probabilities can be read out and inspected by a clinician, which is important for clinical use.

We chose HMM over more complex deep learning architectures like Recurrent Neural Networks (RNNs) or Transformers. Neural networks often reach high predictive performance on medical time series, but they also need large and well-labeled datasets. They are hard to interpret (the so-called black-box problem). They also do not work well with the severe feature missingness which is typical for administrative health data. Recent literature shows that probabilistic approaches like HMMs can give good, interpretable predictions on medical sequences and are also much easier to train and validate~\cite{hmm_medical_time_series_2022}.

\subsubsection{Vector space modeling (TF-IDF and classifiers)}
HMMs model the strict chronological order of events. Another approach is to ignore the order and treat the whole patient history as a ``document'' of medical codes. Then we look only at which codes appear and how often. For this type of representation we used the Term Frequency-Inverse Document Frequency (TF-IDF) statistic.

TF-IDF comes from natural language processing (NLP)~\cite{tfidf_icd_codes_2022}. It is a way to assign weights to terms in a document. In our case a term is a medical procedure code and a document is one patient history. The Term Frequency (TF) part counts how often a given code appears in one patient's history. The Inverse Document Frequency (IDF) part looks at how many patients in the catalog contain this code at all. Codes that appear in the records of almost all patients (for example a routine blood test) get a low IDF weight, so they are effectively down-weighted. Codes that appear only in a small number of patients (for example a specialized oncology biopsy) get a high IDF weight, so they are emphasized. The final TF-IDF weight of a code for a patient is the product of these two.

So every patient history is turned into a fixed-length numeric vector over the full code vocabulary. The original problem, which was to classify a variable-length sequence, becomes a standard classification problem. We trained two classical classifiers on these TF-IDF vectors: Logistic Regression and Random Forest. Logistic Regression models the probability of class membership through a sigmoid function. We used it as a fast baseline. Random Forest is a set of many decision trees. Each tree is trained on a different random subset of the data and a different random subset of features. Each tree votes for a class, and the final prediction is the majority vote across all trees. We used Random Forest to capture non-linear interactions between groups of medical codes and the target disease outcome, which Logistic Regression cannot do directly.


\subsection{Application architecture and technology stack}
\label{subsec:architecture}

The system is designed as a containerized web application. The user interface is separated from data processing and model inference. The architecture uses a client-server model. On the client side there is a Single-Page Application (SPA) frontend. On the server side there is an asynchronous REST API backend, a background task queue for heavy processing, and several specialized databases for different data types. All components are put into Docker containers and started together with Docker Compose. So the same environment is used in development and in production deployment.

\subsubsection{Client-side Single-Page Application (SPA)}
The user interface is implemented as a Single-Page Application using React and TypeScript. In an SPA the browser loads one HTML page and a JavaScript bundle once at the beginning. After that, when the user clicks to another view (for example from the patient catalog to the scoring system), the page is not reloaded. Instead, the application catches the navigation, updates the browser URL, and re-renders only the components that changed. The new data is fetched asynchronously from the backend as JSON~\cite{spa_benefits_digiteum}.

This approach was chosen because it avoids full-page reloads, which is important for multi-step workflows like the sequential data input form in the scoring system. When the page would reload, the user would lose the already-entered codes. Another reason is that React's component-based architecture fits well for data-intensive dashboards. Complex visualizations (for example longitudinal patient heatmaps rendered with Chart.js) can be updated interactively without losing the surrounding application state~\cite{react_data_viz_2025}. TypeScript was used to provide static type checking. So the code is easier to maintain and the data interfaces expected by the REST API are enforced at compile time.

\subsubsection{Backend REST API and asynchronous processing}
The core backend logic is provided by a REST API developed in FastAPI, a Python web framework. FastAPI is built on the ASGI (Asynchronous Server Gateway Interface) standard and natively supports Python's \texttt{async}/\texttt{await} syntax~\cite{fastapi_ml_advantages}. So the API can handle several concurrent I/O-bound requests (for example database queries or fetching patient records) without blocking the main execution thread while waiting for a slow operation~\cite{fastapi_ml_advantages}.

Machine learning inference on sequential health data is computationally intensive. It can take several seconds to complete. So if these predictions were run directly inside the HTTP request-response cycle, the browser connection would time out and the user interface would freeze. For this reason we implemented an asynchronous background job queue with Celery and Redis~\cite{celery_redis_devto}. When a user submits a sequence of medical codes for scoring, the FastAPI backend stores a Queued task state in the database, pushes the task payload to the Redis message broker, and immediately returns a unique Task ID to the client. A separate Celery worker process keeps the machine learning models loaded in memory, listens to the Redis queue, runs the inference, and writes the final prediction result back to the database.

\subsubsection{Polyglot persistence data layer}
The data managed by the application is very diverse. It would be impractical to store everything in a single database engine. So the system uses a polyglot persistence strategy. Different types of data are stored in different database technologies, each one chosen to fit the specific access pattern~\cite{polyglot_dataexpert_2026}.

The persistence layer is split as follows:
\begin{itemize}
    \item \textbf{MongoDB (document database).} It is used to store the unstructured training data catalog. Electronic health records are hierarchical and have variable length. One patient can have five recorded procedures and another fifty. A NoSQL document database like MongoDB stores each patient as a flexible JSON-like document, so the whole longitudinal history can be retrieved as one object without complex relational joins~\cite{mongodb_healthcare}.
    \item \textbf{MySQL (relational database).} It handles structured, transactional data. This includes user accounts, role-based access control (RBAC) permissions, and metadata about the uploaded machine learning models. The relational model enforces a fixed schema and keeps data integrity, which is important for the parts of the system that are security-critical.
    \item \textbf{Elasticsearch (search engine).} It provides fast, fuzzy text search over the dictionary of five-digit medical procedure codes. So the frontend can offer real-time search suggestions when a doctor types a code into the scoring form.
\end{itemize}

\subsubsection{Reverse proxy and entry point}
An Nginx server is used as the reverse proxy and the single entry point for all incoming HTTP traffic. The proxy routes traffic by URL path. Requests for the API (for example \texttt{/api/*}) are forwarded to the FastAPI container. All other requests are routed to the React frontend container. So access control is kept in one place. It also simplifies Cross-Origin Resource Sharing (CORS) configuration. The same Nginx layer is prepared for later production requirements like load balancing or TLS/SSL termination.


\subsection{User-centered design of the application}
\label{subsec:ux_design}

Adoption of machine learning tools in clinical practice is often blocked by poor usability and by bad integration into the existing workflow~\cite{ucd_ml_dashboard_2023}. To reduce these barriers, the development process used a User-Centered Design (UCD) methodology. The UCD work included defining target user personas, doing task analyses, and making iterative wireframes. These steps were done before the full-stack implementation started~\cite{jmir_personas_2022}.

\subsubsection{Target users and personas}
In healthcare software design, personas are used to represent the needs, technical literacy, and typical workflow of the prospective end users~\cite{personas_healthcare_2022}. For this application we defined two main personas, based on the intended use cases:

\begin{enumerate}
    \item \textbf{Clinical specialist / physician.} This user is mainly interested in the scoring system. They need a fast and intuitive interface to enter a patient's medical history (manually or as a dataset) and to immediately see the machine learning prediction. Then they use the prediction to support a clinical decision. So the important things for this persona are clear terminology, fast search for medical codes, and results which are easy to read.
    \item \textbf{Real-world evidence (RWE) analyst / data scientist.} This user is focused on the training data catalog and on model management. They need tools to upload new predictive models, to explore large cohorts of patient data, to filter longitudinal sequences, and to compare ground-truth labels against model predictions (for example with heatmaps). So they can evaluate how the model performs on real-world populations.
\end{enumerate}

Defining these two personas helped us decide how the application architecture and user interface should be organized, because the two groups have different but also partially overlapping requirements.

\subsubsection{Hierarchical Task Analysis (HTA) and UX flow}
To structure the interaction inside the application, we also did a Hierarchical Task Analysis (HTA). HTA is a method that takes a complex goal (for example ``get a disease prediction for a new patient'') and breaks it down into a sequence of smaller sub-tasks and operations~\cite{personas_healthcare_2022}.

For the scoring system, the HTA produced a guided multi-step workflow. The form is not one long complex page. Instead it is split into several sequential steps:
\begin{itemize}
    \item Step 1: Model selection (for example Lung Cancer or Multiple Sclerosis), limited by the user's permissions.
    \item Step 2: Selection of the data input method. The user can enter codes manually through an autocomplete search or upload a CSV dataset in bulk.
    \item Step 3: Data validation and review.
    \item Step 4: Submission to the asynchronous queue and redirection to the history and results view.
\end{itemize}
This flow directly shaped the frontend architecture. It made it necessary to use a centralized state management solution in React, so the form data is kept while the user navigates between steps.

\subsubsection{Visualizing longitudinal data with heatmaps}
For the RWE analyst persona, sequential patient data is hard to read directly as raw arrays of medical codes. So the application uses interactive heatmaps in the catalog detail view. The heatmap puts chronological clinical events on one axis and disease categories on the other. So the user can visually compare the sequence of ground-truth diagnoses written by physicians against the predictions produced by the machine learning model at the same time points. This kind of view makes it easier to spot time-related patterns and to see the moments where the model disagrees with clinical reality.

\subsubsection{Wireframes and usability testing}
After the task analysis, low-fidelity wireframes and interactive mockups were made. They visualize the layout and the interaction patterns. The wireframes were used as a discussion tool between technical constraints and clinical usability~\cite{ucd_ml_dashboard_2023}.

The mockups were tested with proxy users to find friction points. Feedback from these sessions directly influenced several key design choices in the final application. For example:
\begin{itemize}
    \item Implementation of a fuzzy-search autocomplete for the five-digit medical procedure codes. Users found manual code lookup very hard without it.
    \item Addition of visual cues (color-coding and highlights) to distinguish predicted outcomes from ground-truth values in the patient catalog.
    \item Linking of task IDs across the application. So users can easily track the progress of asynchronous prediction tasks on a dedicated history page.
\end{itemize}
The final UI design was then translated into the React component hierarchy described in Section~\ref{subsec:implementation_patterns}.


\subsection{Implementation patterns and state management}
\label{subsec:implementation_patterns}

Developing a responsive single-page application for healthcare data needs careful handling of client-side state, API synchronization, and background task monitoring. For this reason several frontend implementation patterns were used.

\subsubsection{Persistent UI state and URL-based routing}
One of the main problems in single-page applications is keeping the user context across navigation events and page reloads. For the multi-step scoring form, local state management through React hooks was combined with browser storage. So the user's progress (for example manually entered medical codes or uploaded files) is not lost if the browser is accidentally refreshed.

Another pattern we used is URL-based state management. It supports reproducible analyses and shareable views within the training data catalog~\cite{react_url_state_2025}. Instead of keeping filter values (for example a specific searched medical code or pagination offsets) only in isolated component memory, these values are synchronized with the browser's URL query parameters. So when a user searches for a specific code, the code is highlighted in the list view. When the user then clicks through to the patient detail view, the query parameter is passed along, so the highlight context is kept. It also means clinicians can bookmark or share a link to a specific subset of the patient catalog. The exact view is then reconstructed when the link is opened again~\cite{react_url_state_2025}.

\subsubsection{Client-side caching of API responses}
The application often requests large, slow-changing datasets. For example the complete dictionary of five-digit medical procedure codes or chunks of the patient catalog. Fetching these resources over the network on every keystroke or component mount would be slow for the user and would also put unnecessary load on the FastAPI backend.

For this reason the frontend uses an application-level caching strategy~\cite{frontend_caching_2026}. When a user opens the data input form for the first time, the application fetches the medical code dictionary and caches the response. Then later autocomplete queries typed by the user are resolved against this local cache, which gives instant suggestions. In the same way, catalog pagination uses the cache to serve previously visited pages without a redundant database query. So the user interface stays responsive, which is important for clinical workflow.

\subsubsection{Asynchronous task polling and status propagation}
As described in Section~\ref{subsec:architecture}, machine learning inference is offloaded to a Celery background worker. The REST API returns a response immediately after the task is put into the queue. So the frontend client has to actively track the task progress to tell the user when it is finished.

This is done with an HTTP polling pattern~\cite{aws_async_api_2021}. When the scoring form submits a sequence of codes, the FastAPI server responds with a unique Task ID. The React client then enters a polling state. It sends periodic GET requests to a dedicated status endpoint using this Task ID. When the task status changes from Pending to Success, the client fetches the final prediction.

Another step that improves the user experience is that this asynchronous tracking is integrated into the routing flow of the application. After submission, the Task ID is propagated through URL parameters to the History view. So the user can navigate away from the scoring form, open the full prediction history, and see their most recent task highlighted and updating in real time as the background worker processes it~\cite{aws_async_api_2021}.